\documentclass[11pt]{article}
\usepackage[margin=0.9in]{geometry}
\usepackage{amsmath,amssymb,amsthm,booktabs,longtable,array}
\usepackage[T1]{fontenc}
\usepackage{lmodern}
\usepackage[colorlinks=true,linkcolor=blue,urlcolor=blue]{hyperref}
\usepackage{microtype}
\newtheorem{proposition}{Proposition}[section]
\newtheorem{lemma}[proposition]{Lemma}
\newtheorem{theorem}[proposition]{Conditional theorem}
\theoremstyle{remark}
\newtheorem*{status}{Status}
\newcommand{\R}{\mathbb R}
\newcommand{\E}{\mathbb E}
\newcommand{\tr}{\operatorname{tr}}
\newcommand{\supp}{\operatorname{supp}}
\newcommand{\cum}{\operatorname{cum}}
\newcommand{\Rea}{\operatorname{Re}}
\newcommand{\Ima}{\operatorname{Im}}
\newcommand{\norm}[1]{\left\|#1\right\|}
\newcommand{\abs}[1]{\left|#1\right|}
\title{Toward a four-dimensional Yang--Mills construction\\
\large Notebook 01: proof obligations and verified reductions}
\author{Research working document}
\date{17 October 2026}
\setlength{\emergencystretch}{5em}
\begin{document}
\maketitle

\begin{abstract}
The goal is a complete existence argument for a nontrivial
four-dimensional Yang--Mills quantum theory. A first intermediate
target is a nonzero, fully Euclidean-covariant continuum observable
sector obtained from one subsequence of four-dimensional $SU(2)$
Wilson theories. Cutoff uniqueness is unnecessary for that target.
This notebook gives explicit conditional reductions for compactness, factorial
correlation bounds on joint flat tests, rotation restoration, and
positive-energy nontriviality. Fixed-side patch grids repair the all-order
geometry conditionally on uniform translated-grid source estimates.
An independent free-lattice calculation reduces one nontriviality test
to fixed relative kernel accuracy.
It also gives counterexamples to several tempting shortcuts.
The model-specific source estimates, the interacting rotation input, and
the passage from one scalar sector to Yang--Mills remain open.
Companion Notebook 02 now gives a fixed-coupling proof preserving diagonal
reflection positivity in the original coordinate-torus infinite-volume
limit, with the original odd-$L$ regulator. It supplies no uniform
ultraviolet estimate.
Companion Notebook 03 proves the local current commutator with explicit
constants and a convergent Gaussian sharp-cutoff cluster expansion.
Its remaining interfaces are the combined current at complex decoupling
parameters and strict locality in the background variables.
For the Gaussian cutoff part of the quadratic coefficient, it also
proves a direct Hessian bound and volume-limit criterion that avoid
strict background-support activities, conditional on actual uniform
source-response estimates and their identification.
Notebook 04 establishes strict background locality of the Gaussian
cutoff logarithm under a uniformly small local precision perturbation.
Matching that hypothesis, its domain and its symmetries to the actual
Yang--Mills Gaussian is still required.
Notebook 05 proves a quantitative flat axial-to-exponential gauge bridge
with the actual averaged-path penalty, and an interacting cutoff theorem
for disjoint fixed-dimensional cells. It also identifies the correct
physical-field radius estimate and the full exponent's linear Jacobian
term. Notebook 18 extends that integration theorem to overlapping
supports under an explicit weighted local coefficient norm.
Notebook 06 supplies direct scale-uniform flat quadratic bounds through
curvature averaging, a finite-block axial gap, and an explicit compact
bounded-energy right inverse. Their homogeneous identification with the
constrained-minimizer Hessian now follows from the nonlinear construction
of Notebook 15 and the explicit match in Notebook 17.
Notebook 07 proves exponential locality of this flat
action and retained covariance, and their refinement limit at rate
$O(M^{-2})$ in an exponentially weighted kernel norm, uniformly in
compatible periodic volume. Its explicit sinc-alias limit is a free
Gaussian reference. Notebook 19 proves the physical homogeneous
complex-background perturbation in the exact source coordinates;
the interacting continuum limit remains separate.
Notebook 08 constructs the corresponding fine minimizing extension
in the deposited homogeneous flat source gauge. It proves uniform
exponential locality for the potential, first fine derivatives,
curvature, and componentwise Laplacian, with physical H\"older
bounds below the endpoint for the first derivatives. These local
flat response maps are inputs to a further perturbation argument;
they do not establish that argument on a nonflat current domain.
Notebook 09 proves locality of the full flat fine inverse, with
two total fine derivatives in physical cell Hilbert norm, and a
refinement-uniform complex neighborhood for the deposited auxiliary
current. Its coarse precision has a quadratic zero-current remainder
and all-order localized derivative bounds.
Notebook 10 identifies a local projection onto feasible curvature
variations and gives a full weighted expansion for a small curvature
metric perturbation at fixed averaging constraint.
Notebook 11 bounds the separate nonlinear constraint-current term by
composing the corrected fixed-$L$ quadratic jets and summing their
physical scale factors. It establishes a common weighted Hilbert
domain for the full flat current family along $M=L^j$, including
factorial current derivatives. Nonflat and retained-history domains
remain separate.
Notebook 12 strengthens the flat inverse to arbitrary-forcing
pointwise $C^{1,\alpha}$ bounds for $\alpha<1$, and transfers the
full flat current to the homogeneous source norms $N,N_\beta$
on a smaller uniform complex ball. Its current requires only a
supremum bound. It also proves the one-form right-divergence
members and the separate constrained current member. For every
fixed exponent in $(0,1)$, an improved heat-integral argument
removes the H\"older loss. A scalar proof requires an explicit
uniform positive penalty interval; the baseline Laplacian estimate
is not asserted for nonzero current.
Notebook 13 proves a uniform analytic ball for the complete
fixed-reference corrected nonlinear average and an exact local
chart of its constraint. The chart has a uniformly local right
inverse for its ordinary differential.
Notebook 14 computes the exact change to the moving-reference
derivative, including its endpoint correction. A cumulative gauge
relation controls the ordinary reference trajectory and the actual
bond and scalar averages. A Wilson plaquette calculation then gives
a uniformly local energy inverse on a small global physical $C^1$
ball, including the moving gauge projection. Its bare Wilson action,
physical pairing and left-exponential Hessian chart match the
deposited conventions in dimension four. Original-paper verification,
the precise graded-curl and effective-Hessian identifications,
nonflat pointwise/current estimates and retained-history domains
still require their separate checks.
Notebook 15 closes a nonlinear classical minimum construction in
the homogeneous small-global-field class. An equal-exponent
constrained inverse, the exact Wilson force, and a two-Hilbert
constraint-curvature estimate give a uniform analytic fixed-slice
minimum. A genuine constraint-preserving gauge flow and a separate
ordinary-average transfer produce a minimum modulo gauge for the
ordinary homogeneous constraint. Its radii and physical
$C^{1,\alpha}$ bounds are uniform in refinement and volume.
This does not supply mixed histories, large-field integration,
or the interacting continuum observable hierarchy.
Notebook 16 proves exponential locality of both selected
minimizer representatives and their multiplier, at small
nonconstant backgrounds, with one weighted direction at every
analytic order. The exact action-envelope identity gives a local
nonlinear classical coarse Hessian, its small weighted perturbation
from $q_M$, and a classical coarse gauge identity. A positive
auxiliary coarse penalty gives a stable local inverse. Quantum
fluctuation integration and multi-insertion tree estimates are
not supplied by those classical response bounds.
An exact constrained-Gaussian trace calculation in Notebook 16
also shows that its physical curl variance grows as $\eta^{-4}$.
The probability of a fixed smooth ball for a deterministic
background plus $gZ$ is at most $C r\eta^2/g$.
Notebook 17 supplies that rescaling and joint local-tail analysis.
It proves real correlated bad-block bounds and a complex marginal
estimate without a total-volume factor, and compares the actual
first-step reduced Wilson Hessian with the flat energy, including
constraint curvature. It also identifies the homogeneous flat
current-step covariance uniformly over all earlier averaging depths.
Notebook 18 proves the overlapping-interaction cutoff expansion and
its quantitative comparison with the Gaussian logarithm. Under a
small local precision decomposition it also supplies strictly local
background activities. Notebook 19 constructs the physical stationary
background and its exact pivot-coordinate Gaussian, including both
action curvature terms, uniformly in preceding classical depth.
It also bounds the local Jacobian, Haar and gauge interactions.
Notebook 20 closes the actual bare Wilson small-field integration:
the action and precision have explicit local decompositions, the
retained cutoff is fixed, and the logarithm has a volume-uniform
connected expansion. Its background-only delta and Gaussian
determinants are included separately. This support statement is
in independent current backgrounds, before composition with the
next-coarse stationary field. Notebook 21 proves strict support
localization of the homogeneous classical minimizer and action
uniformly in earlier refinement. A rooted support contraction,
with arbitrary fixed cardinality reserve, gives local scalar
action coefficients and individually symmetric precision pieces.
The physical chart substitution then closes one small-field
integration of that classical reference action at every earlier
classical depth. Notebook 22 constructs the actual stationary
current background as a strictly supported analytic function
of next-coarse data and transports the generated activities
through it. That bounded step has explicit analytic-radius
and support-weight losses; it does not give repeated-step
closure. Notebook 23 proves local gauge covariance of the
complete small-field correction, including the cutoff,
Haar and pivot factors, and Gaussian determinant. Its total
linear jet vanishes and its quadratic jet obeys the finite-regulator
Ward identity. Notebook 24 bounds the inherited fluctuation
increment by $O(gR)$ divided by the available analytic margin
and uses the
exact pivot support geometry to give a sufficient next-step
interaction-norm inequality. The resulting logarithm comparison
is proved under that inequality. Notebook 25 gives a smaller
pre-root-gauge representative with an inverse-ratio link-angle
bound and a strict support series, uniform through iterated
classical blocking. It extends the one-step composition germ of
a gauge-invariant total functional where that smaller representative
is admissible.
Notebook 26 supplies independent auxiliary integration cells
that satisfy the support-weight test: a fixed positive
connected-carrier exponent is recovered through inherited
insertion, Gaussian integration, all normalization terms and
stationary transport. The physical blocking and cutoff stay
fixed. This is a bounded step on a smaller domain, with explicit
noncontractive constants. Notebook 27 constructs the actual
composite ordinary blocking map, its local pivot chart and
retained Gaussian, and the precisely normalized compact-gauge
integral. Notebook 28 composes around a local core of the
pre-gauge representative, expanding only its small distant tail.
Classical support substitution before integration then gives
one small-field update on the same input and output analytic
ball, with a fixed positive carrier exponent and a nonvanishing
Gaussian cutoff integral. It also proves an inverse-ratio
background contraction for coefficientwise Taylor remainders
of field degree at least five, and a contractive quantum
response to input differences with matching low jets.
The update can regenerate low jets and the general activity
bound can grow. Notebook 29 proves the actual local spatial
symmetries and classifies the total quadratic, cubic and
quartic marginal vertices through their nonlinear Ward
identities. One normalized gauge-invariant classical action
supplies their common subtraction. Complete first-gradient
H\"older support families and an exact lattice difference
factorization give a connected-norm spatial gain
\(C(J)L^{-\alpha}\) under the stated infinite-kernel and
coherent-family hypotheses. A classical semigroup modulo
gauge permits \(L=R^q\) to be chosen after fixed \(R,J\).
Joint compactness realizes the coherent infinite interactions
at fixed parameters. A complete quotient by equality of
normalized totals makes classical representation changes
and bounded full-reference extraction precise; the resulting
classical contraction retains its joint-map hypotheses.
Notebook 30 closes the universal gauge identity for that fixed
classical system and constructs a norm-preserving finite
periodization of invariant interaction classes. Strict Gaussian
local continuity, with explicit carrier-size and expansion-order
tails, defines the actual homogeneous local quantum map on the
quotient. It identifies varying finite inputs and every fixed
finite history whose intermediate finite families meet the
uniform budgets. Exact transport of the actual classical
reference makes intermediate extractions telescope. A finite
block of projected quantum steps has an \(O(b^2)\) stable amplitude
envelope while its couplings remain uniformly weak. A fixed outer
domain supplies the matching \(O(b^2)\) zero-input force, and
each inverse-coupling increment has a bounded absolute size.
Notebook 31 closes the finite representative-budget interface.
Finite Ward repair and periodic localization retain an explicit
wrapping error at the same field radius and carrier exponent.
Exact finite quotient descent and equivalent small representatives
then give uniform local induction through every permitted finite
history while its couplings remain weak. All raw constants are
compensated in the vacuum factor. A fixed-threshold parity
argument isolates the pivot and Gaussian determinant as the
leading zero-input correction. On the jointly compatible infinite
interaction space, a quadratic-only spatial estimate at fixed
radius proves exact classical preservation of the Maxwell
coefficient. The leading quantum coefficient and its drift sign
still require evaluation; the infinite identity does not discard
finite cutoff-extraction errors.
Notebook 32 proves qualitative uniform disappearance of that
finite classical leakage along cofinal nested periods and joint
continuity of the finite quantum map in its class and positive
real coupling. It derives the actual leading determinant's
convergent Hessian trace and a finite composition identity from
one Haar pushforward evaluated directly and sequentially.
Together with classical Maxwell preservation this gives additive
infinite leading coefficients. An exact gauge/coarea identity
retains vector, based-ghost and coarse-constraint determinants;
a Rayleigh test rules out a uniform physical mass bound on the
based ghost alone. Identifying the coarse constraint covariance
with a shifted actual effective-Hessian inverse gives a uniform
bound on that coarse determinant's Maxwell contribution for a
fixed sufficiently large shift. A positive block margin uniform over all
varying-depth structural limits would transfer to actual finite
regulators by qualitative compactness, without a rate.
A separate shooting theorem tunes
arbitrarily long finite local histories to a fixed weak endpoint
under an explicitly unproved positive averaged drift hypothesis
for the actual finite coefficients.
Notebook 33 identifies individually gauge-invariant diagonal
determinant densities with the total Maxwell moment. Polynomial
stationary jets on a constant Cartan-curvature probe remove the
multiplier-curvature contribution through the required orders,
including for the effective classical actions. A specified massive
continuum proper-time reference has an explicit positive logarithm
and bounded finite part in the same color and extraction convention.
An exact heat integral for the actual lattice coefficient isolates
sufficient estimates on three heat-time intervals. A closed-walk
expansion bounds the actual based-ghost short-time contribution
uniformly in endpoint and preceding depth. The vector contribution,
the other intervals, and hence the sign for the actual family
remain open.
Notebook 34 computes the entire actual based-ghost coefficient.
A returning-walk area formula, residue counting, the sharp
point-pinned gap and a fourth-moment tail estimate give
\(C_H\log K/(12\pi^2)+O(1)\), uniformly in preceding depth.
Its extra logarithm on heat times \(K^2\) to \(K^4\) proves
that separated long-time ghost control at the usual block
scale is impossible. The notebook also bounds the full
actual vector short-time contribution, preserving the
composite constraint. Thus the ultraviolet heat interval
is closed; the intermediate and longer-time coupled
comparison remains.
An exact finite physical-quotient identity removes
periodic pinning while retaining the fine/coarse gauge-orbit
density ratio. The residual physical precision has a uniform
flat gap \(cK^{-2}\). A contour-gauge counterexample shows
why the unprojected contour matrix cannot be assigned that
gap. A paired scalar reference has exact infrared cancellation;
its identification with the actual finite factors and
uniform nonflat projected estimates remain open.
Notebook 35 resolves the fixed-endpoint infrared passage by a
new exact local gauge choice. A periodic flat-holonomy calculation
first identifies a noncommuting scalar limit and its forced
physical compensation. Moving-reference endpoint tangents then
replace point pinning by zero covariant block means, with exact
constant row Gram matrix. Finite Gaussian and Schur identities,
continued jointly through flat stabilizers, express the actual
determinant as a gauge-covariant scalar/vector pair. Both factors
have flat gaps of order \(K^{-2}\); a block-mean harmonic lift
and a local tent lift prove full vector coercivity and Laplacian
domination. Uniform heat and inverse-power bounds follow.
Their local diagonal densities identify the actual Maxwell
coefficient at every fixed endpoint. The large-endpoint
curvature comparison and actual positive drift remain open.
Notebook 36 factors that actual determinant into a soft fine
scalar/vector pair and two coarse determinants. Both coarse
Maxwell contributions are uniformly bounded as the endpoint
grows. The proof controls physical source kernels without a
fine-fiber dimension loss. A local gauge argument and an
ordered Laurent factorization give uniform relative vertices
for the actual effective-action Hessian; the vector matching
uses its established on-shell affine jets. The remaining fine
pair retains the effective action and its constraint terms.
Both ends of its heat integral are now uniformly bounded:
the current-time interval below one and the physical infrared
tail above \(K^2\). A local cubic-resolvent argument proves
the latter without cyclically moving a block projector.
The actual Maxwell logarithm is thereby reduced, up to a
uniformly bounded term, to heat times \(1<s<K^2\).
Notebook 37 closes this leading-determinant comparison under the
explicit classical structural and local quantum inputs of the series.
Source-localized sectorial resolvents remove the soft penalties on
the whole band up to physical time one with an \(O(u)\) heat error.
The total nonlinear Ward jets through degree four then compare the
actual effective-action base with the Wilson base with error \(O(s^{-1})\),
retaining the left-relative chart-current terms. A direct
Wilson-to-continuum comparison has error \(O(s^{-1/2})\), with its
color and determinant normalization checked both by the magnetic
heat kernel and by the transverse two-point diagrams. Consequently
\(\beta(\mathcal D_{M,K})=22\alpha_g\log K+O(1)\), uniformly in old
depth and every applicable joint structural system. This supplies
the positive margin required by finite moment transfer and local
endpoint shooting. Arbitrarily long finite weak homogeneous local
histories can be tuned to one fixed weak endpoint, with bare coupling
of order the inverse square root of their length. These are research
proofs from the stated upstream inputs, not external certification
of the complete series or a continuum Yang--Mills existence theorem.
Notebook 38 supplies exact positive good/remainder operators
for a specified measurable completion of the local coarse map.
Their iterates preserve the full Wilson mass and every history;
local-germ agreement requires literal real-window containment
and a possibly finite-volume-dependent orbit-overlap shrink.
Separately, exact site-plane gluing gives uniform-volume
bad-plaquette Peierls bounds under the normalized bare Wilson
measure. A fixed-box conditional theorem includes its full
compact normalization, and a boundary-only gauge joins these
results into joint rarity of explicit local gauge/minimizer
window defects. These box windows are not yet the constrained
windows of the multiscale map. Their identification, normalized
large-component replacements, complete-density induction, and
the continuum observable and reconstruction arguments below
remain. Notebook 39 completes a restricted normalized replacement
on arbitrary finite families of separated fixed boxes in the bare
Wilson measure. It preserves the exterior marginal and raw mass,
leaves bad-boundary laws unchanged, and gives exact sparse failure
and coupling formulas. Its local total-variation, entropy, cluster
and bounded-source cumulant comparisons are independent of the
number of distant boxes. This finishes a finite local measure
operation; it does not identify it with the constrained multiscale
step or prove continuum existence.
\end{abstract}

\noindent\textbf{Evidence boundary.}
The repository was examined at commit
\href{https://github.com/ethanwsussman/Yang--Mills-LLM-experiment/tree/e392ef45013a581664bab76b90867170b30fe10e}
{\texttt{e392ef4501\allowbreak 3a581664ba\allowbreak b76b908671\allowbreak 70b30fe10e}}.
The claims below concern the deposited material; unposted companion proofs
have not been assessed. Elementary lemmas are proved here. Their applicability
to the proposed Yang--Mills limit is always a separate assertion.
No novelty or independent certification of the entire manuscript is claimed.

\section{The target and the actual missing pieces}

We should distinguish three milestones.
\begin{enumerate}
\item \textbf{Continuum scalar sector.}
Construct a consistent family of separated Schwinger functions for a
renormalized action-density observable, with reflection positivity, full
Euclidean covariance, sufficient all-order distributional bounds, and
nonzero positive-energy content.
\item \textbf{Yang--Mills theory.}
Extend to an adequate family of gauge-invariant curvature polynomials and
their derivatives, including mixed correlations and the correct tensor
transformation laws. Identify their normalization, short-distance behavior,
and identities along the Wilson continuum trajectory. Establish the
appropriate vacuum-sector and reconstruction properties.
\item \textbf{The full existence and mass-gap problem.}
In addition, prove a strictly positive gap and address every compact simple
gauge group. This is stronger than the first two milestones.
\end{enumerate}
The official problem statement explicitly asks for gauge-invariant local
fields and appropriate short-distance behavior, together with the quantum
field axioms and a mass gap \cite{clay}. One scalar called
$\tr F^2$ does not automatically provide that observable family.

\begin{center}
\small
\begin{tabular}{p{0.23\linewidth}p{0.32\linewidth}p{0.36\linewidth}}
\toprule
Obligation & Present status & Best next proof target\\
\midrule
Uniform local correlations &
Conditional source/RG inputs &
One fixed physical stopping scale, a common vacuum normalization, and
uniform complex-source bounds.\\[3pt]
Subsequential existence &
Follows conditionally from local bounds &
Use a countable diagonal extraction; defer cutoff uniqueness.\\[3pt]
All-order reconstruction bounds &
Conditional factorial and joint-test reduction below &
Verify the translated-grid source estimates and anchored polymer
bounds that feed the reduction.\\[3pt]
Full rotations &
The supplied two-form theorem does not directly cover $\tr F^2$ &
Pass the two-frame positivity proved in Notebook 02 to the renormalized
hierarchy, and prove an angular growth exponent strictly below four.\\[3pt]
Nontrivial dynamics &
The proposed lower witness is conditional &
Compare two explicit reflected correlators at the same scale.\\[3pt]
Yang--Mills identification and vacuum &
Not supplied by a scalar witness &
Construct mixed gauge-invariant observables and justify phase selection
in the actual observable framework.\\
\bottomrule
\end{tabular}
\end{center}

An ordinary positive probability law for the coincident composite Euclidean
field is \emph{not} a prerequisite we should impose without reason.
The draft itself uses test functions flat on coincidence diagonals.
An appropriate Osterwalder--Schrader reconstruction on that domain is a
possible route, provided its quantitative growth assumptions are proved
\cite{os2,seiler}. Conversely, fixed-order separated correlations alone do
not establish those assumptions.

\section{A smaller source-stability interface}

Write the finite-lattice generating function as
\[
 Z_{a,\Lambda}(w)=\int e^{-x_a A(U)+\mathcal O_a(w)}\,dU,
 \qquad \widehat Z_{a,\Lambda}(w)=e^{-C_{a,\Lambda}}Z_{a,\Lambda}(w).
\]
Here $a$ is the ultraviolet spacing and $C_{a,\Lambda}$ is a real,
source-independent vacuum normalization. It cancels from all ratios.
It is essential: the unnormalized bare free energy has a cost proportional
to the number of \emph{bare} sites. We require estimates proportional to
the number of sites at one \emph{fixed physical coarse spacing} $h$.
Put $V=|\Lambda|/h^4$.

\begin{lemma}[Elementary denominator floor]
Suppose an exact coarse integral identity gives
\[
\widehat Z(0)=\int \rho(V_c)\,dV_c,\qquad
\rho(V_c)\geq\chi(V_c)e^{-x_*A_h(V_c)-E_-V},
\quad 0<x_*\leq X_*.
\]
Suppose $\chi=1$ when every coarse link lies in a fixed group ball
$B_\delta$ about the identity, of normalized Haar measure $H_\delta>0$.
If $A_h\leq 6C_A\delta^2V$ there, then
\[
\widehat Z(0)\geq e^{-C_-V},\qquad
C_-=E_-+6C_AX_*\delta^2+4\log H_\delta^{-1}.
\]
All constants are independent of $a$ and $\Lambda$.
\end{lemma}
\begin{proof}
Restrict the integral to that product of Haar balls. There are at most
$4V$ links and $6V$ plaquettes. The restricted measure is at least
$H_\delta^{4V}$, and the asserted lower bound follows.
\end{proof}

The small-field characteristic function $\chi$ is a real hypothesis:
in the RG construction it involves minimizing backgrounds, not merely
the visible coarse plaquettes. Its uniform window must be verified.
This elementary integration does not prove that window.

\subsection{The fixed-cell reduction}

Tile the coarse torus by $M$ congruent reflection cells, each of volume $v$.
Let $w_z=\sum_{i=1}^n z_i f_i$ be supported inside one cell, with a positive
buffer from its faces. The $f_i$ are real, with separated supports.
Write $\nu_i=\norm{f_i}_\sharp$ and
$D_r=\{z:\sum_i|z_i|\nu_i<r\}$.
Here $\norm f_\sharp=\max(\norm f_\infty,\operatorname{Lip}f)$
in the fixed coarse physical units.
Let $w_z^{\rm ch}$ be its chessboard replication, conjugating coefficients
on odd reflection copies.

Assume the following quantitative inputs.
\begin{enumerate}
\item The original Wilson measure gives
\[
 \abs{Z(w_z)/Z(0)}
 \leq [Z(w_z^{\rm ch})/Z(0)]^{1/M},
 \qquad Z(w_z^{\rm ch})/Z(0)\geq0.
\]
\item A local analytic source counterterm $E_a$, with $E_a(0)=0$, obeys
\[
 \abs{e^{-E_a(w_z^{\rm ch})}\widehat Z(w_z^{\rm ch})}
 \leq e^{C_+Mv}.
\]
\item Split $E_a$ into cell-interior terms and face terms, each normalized
to vanish at zero source. The latter obey
$|E_a^\partial(w_z)|\leq C_F$ and
$M^{-1}|E_a^\partial(w_z^{\rm ch})|\leq C_F$.
There is a holomorphic recentering $E_a^\#(z)$, normalized by
$E_a^\#(0)=0$, with
\[
 \Rea E_a^\#(z)
 =M^{-1}\Rea E_a^{\rm bulk}(w_z^{\rm ch})
   +\Rea E_a^\partial(w_z),
 \qquad
 |\partial_1\cdots\partial_n E_a^\#(0)|
 \leq B_n\prod_i\nu_i.
\]
One obtains the recentering by averaging the transported bulk terms,
conjugating outside the terms from odd cells as well as inside their
coefficients. This double conjugation is holomorphic.
\end{enumerate}

\begin{theorem}[Uniform fixed-order correlations]
Under these inputs and the denominator floor, set
\[
h_a(z)=e^{-E_a^\#(z)}Z(w_z)/Z(0),\qquad
Q=e^{(C_++C_-)v+2C_F},\qquad r_0=\frac{r}{2Q+1}.
\]
Then $h_a(0)=1$, $|h_a|\leq Q$ on $D_r$, and $|h_a-1|\leq1/2$ on $D_{r_0}$.
In particular,
\[
 \abs{\cum(\mathcal O_a(f_1),\ldots,\mathcal O_a(f_n))}
 \leq
 \left[B_n+\log2
 \left(\frac{n(2Q+1)}{r}\right)^n\right]\prod_i\nu_i.
\]
\end{theorem}
\begin{proof}
Combine the chessboard inequality with the source upper bound and
denominator floor. The bulk real parts cancel exactly against $E_a^\#$;
the two face errors cost at most $2C_F$. This proves $|h_a|\leq Q$.
On the complex line from zero through $z$, the Cauchy coefficient bound
gives $|h_a(z)-1|\leq Qt/(1-t)$ for
$t=r^{-1}\sum_i|z_i|\nu_i$. Hence a local logarithm has modulus at most
$\log2$ on $D_{r_0}$. Cauchy on radii $r_0/(n\nu_i)$, followed by the
mixed derivative of $E_a^\#$, proves the result.
\end{proof}

The constants are independent of cutoff and full volume. A countable dense
family of separated test tuples therefore admits a common convergent
subsequence. Uniform multilinear bounds extend the limit to arbitrary
separated tuples. Use cumulants of orders at least two to form moments
of the \emph{original} lattice observable centered by its real
expectation. Reflection positivity passes through finite matrices of
those moments. It is not asserted for cumulants or for the
nonlinearly recentered generating function $h_a$.
Cutoff uniqueness is not used.

One cell enclosing all insertions does \emph{not} provide bounds uniform
when two groups move arbitrarily far apart. For a contraction time
semigroup one needs a per-ball version: finitely many occupied cells of
fixed side, with the same face buffer, yielding a cost proportional to
their number rather than the diameter of their union.

\subsection{A checkable locality condition}

At physical scales $\ell_j=L^{-j}$, one sufficient counterterm estimate is
\[
\sum_{\substack{X^+\cap B_R\ne\varnothing\\ \operatorname{diam}X^+\geq t}}
 \sup_w |e_{a,j}(X;w)|
\leq A(R+\ell_j)^4\ell_j^{-4}(1+j)^p e^{-ct/\ell_j}.
\tag{2.1}
\]
Each term must depend only on the source on $X^+$.
The scale sum satisfies
\[
H_p(t):=\sum_{j\geq0}(1+j)^p L^{4j}e^{-ctL^j}
\leq C(1+t^{-4})[1+\log^+(1/t)]^p.
\]
Split at $L^j\simeq1/t$ and sum the geometric tail to prove this.
A mixed derivative involving separated supports kills every term that
does not meet them all. Surviving terms have diameter at least their
separation. The same observation controls face errors using the buffer.
Short local divergent terms are allowed; they must not be confused with
the long terms whose sum is required to converge.

\begin{status}
This isolates a smaller verification task; it does not prove the source
upper estimate or (2.1). Unsourced ultraviolet stability cannot be
substituted for these source-dependent estimates. The deposited account
already uses recentering and chessboard arguments; the reduction here
clarifies the necessary interface and removes unnecessary uniqueness
requirements.
\end{status}

\subsection{The analytic estimate that the source restriction does not remove}

The main source bound is supplied by the proposed correlation theorem in
Raw II, lines 2791--2979, and its last-lattice argument at
59645--59855. It is not the later ``source lemma'' at
145173--145219, which is an additional input to the lower two-point
witness. The bare density at lines 737--745 already includes a
source-independent vacuum subtraction; that convention must be retained.

For a source $f$ buffered from cell faces, let $F_e$ denote its reflected
copies and let $F_o$ additionally change sign with reflection parity.
Reflection and vanishing near the faces give
$\norm{F_e}_\sharp,\norm{F_o}_\sharp\leq\norm f_\sharp$:
split a line segment into its portions in individual cells and add their
Lipschitz bounds. Thus
\[
 (zf)^{\rm ch}=(\Rea z)F_e+i(\Ima z)F_o
\]
lies in the advertised two-profile source domain when
$\sqrt2|z|\norm f_\sharp<r$. This proves source-domain membership,
not source stability. The independent complex decoupling variables
used to localize the RG terms remain present even at zero source.

\textbf{Dependency correction.}
The three localized links at Raw II, lines 48058--48320, describe site
(L) in one part of the account. It has not been established that this
reading is a required bottleneck of the active source proof. Its main
induction consumers use
a \emph{global four-layer chain}. Raw I, lines 73107--73123 and
73545--73634, identifies that chain, explicitly excludes the alternative
localized reading, and makes global carrier containment the base of
the holomorphy argument. Its layers have both current and block-data
inputs. The next model obligation is to verify this actual layer's
joint domains, four-stage membership, and inherited or arrested slots.
The following sufficient criterion is abstract: no identification of
the global layers with these fixed-point equations is presumed.

\begin{lemma}[Four-stage sufficient domain criterion]
For field slots $j$, equip the full analytic chart norms with radii
\[
 \delta_j=A_1g_j(\log g_j^{-2})^{p_0},\qquad
 \alpha_j=C_*g_j(\log g_j^{-2})^q,\quad q>p_0,
 \qquad \norm u_\delta=\max_j\norm{u_j}_j/\delta_j.
\]
The norms must include all derivatives and other seminorms defining the
charts. Put $\epsilon=\max_j\delta_j$. Suppose four response maps
solve, successively,
\[
 H_i=G_i(\sigma_i)\xi_i+
       G_i(\sigma_i)\mathcal N_i(H_i;\sigma_i,\psi_i),
       \qquad i=1,2,3,4,
\]
where each output becomes the next input. Assume jointly holomorphic
coefficients on the required charts and decoupling polydisc, with
\[
 \norm{G_i}_{\delta\to\delta}\leq D,\quad D\geq1,
 \qquad \mathcal N_i(0)=0,
\]
\[
 \norm{\mathcal N_i(u)-\mathcal N_i(v)}_\delta
 \leq A\epsilon(\norm u_\delta+\norm v_\delta)\norm{u-v}_\delta.
\]
Assume these bounds on the balls below, and that the background occupies
at most half of every analytic radius. If
\[
 64AD^5\epsilon\leq1,\qquad
 32D^4\delta_j\leq\alpha_j\quad\hbox{for every }j,
\]
then inputs $\norm{\xi_1}_\delta\leq1$ give jointly holomorphic outputs
of norm at most $2D,(2D)^2,(2D)^3,(2D)^4$, respectively. All outputs fit in
the available analytic radii, uniformly in the number of slots.
\end{lemma}
\begin{proof}
For input norm $s$, use the ball of radius $R=2Ds$.
The fixed-point map has Lipschitz constant at most
$2DA\epsilon R=4AD^2\epsilon s$. Even at the largest input
$s=(2D)^3$, this is at most $1/2$.
Its linear part is bounded by $R/2$, and its nonlinear part by
$DA\epsilon R^2\leq R/4$. It maps the ball into itself and contracts.
Uniform Picard iteration gives holomorphy. Repeat four times and use
$(2D)^4\delta_j\leq\alpha_j/2$ for domain containment.
\end{proof}

At the same scale, a fixed small coupling ceiling $\gamma$ supplies
the required radius margin as soon as
\[
 (\log\gamma^{-2})^{q-p_0}\geq32D^4A_1/C_*.
\]
It also supplies small $\epsilon$; no finite bound on the number of UV
steps is needed. The difficult estimate is the \emph{weighted} operator
norm, for example the sufficient row bound
\[
 \sup_j\sum_\ell\norm{G_{j\ell}(\sigma)}
             \frac{\delta_\ell}{\delta_j}\leq D.
\]
Ordinary spatial decay does not establish this when a current field is
substituted into an arbitrarily old slot with a smaller analytic radius.
The two-slot map $T(u_j,u_k)=(u_k,0)$ has unweighted norm one, yet
$T(0,\delta_k)$ escapes a radius $\alpha_j\to0$.
This is a diagnostic for the proposed estimate, not a counterexample
to the actual localization maps.

There is a more specific positive result for one response chart actually
displayed in Raw I. Its equations, at lines 14376--14437, have the form
\[
 A_0=T[\Theta_JH_0B-V'^J(A_0+H_0B)],\quad A=A_0+H_0B,
 \quad\mathcal H^J(B)=A-H^JD^J(A).
\]
Direct differentiation gives
\[
 D_B\mathcal H^J=(I-H^JD'(A))
 (I+TV''^J(A))^{-1}(I+T\Theta_J)H_0.
\]
The normalization includes
$H_0=G_0Q^*(QG_0Q^*)^{-1}\ell^{-1}$, where $\ell$ is the local
layer-scale operator. It cannot be dropped.
The coupling-flow inequalities in Raw I, lines 10234--10245, imply
\[
 \delta_k/\delta_j\leq C_w\sqrt{1+|k-j|}
\]
once the fixed coupling ceiling is sufficiently small. For example,
$x_j=g_j^{-2}\geq\gamma^{-2}$ and
$-c_5\leq x_j-x_k\leq B^\sharp(k-j)$ for $j\leq k$ give this by
bounding the square-root and logarithmic factors separately, with
$B^\sharp\gamma^2\leq1$ and $\log\gamma^{-2}\geq2p_0$.

For this chart the multiscale metric obeys
$d(y,y')\geq W(|j-j'|-1)_+$, as recorded at Raw I,
lines 13471--13483. Hence, if $\mu W\geq1/2$,
\[
 e^{-\mu d(y,y')}\delta_{j'}/\delta_j
 \leq\sqrt2 C_w e^{-\mu d(y,y')/2}.
\]
The weighted row bound therefore follows from the claimed response
kernel estimate and block summability in that metric. In this
geometry, scale-age weights are not an additional independent obstacle.
The response kernel estimate itself and its full-polydisc antecedents
still require verification. The $B$ derivative alone does not cover
the current input of the selected global layer. An accessible copy of
the original CMP 122-II paper is now available; its formulas (1.51)
and (1.54)--(1.55), on printed pages 370--371, are relevant to that
global route. The alternative localized formulas (1.61)--(1.64) must
not be substituted for it merely because some operator letters agree.
Source extraction of several displayed formulas is incomplete, so
equation-by-equation identification remains subject to that evidence
limit. See \cite{balaban122b}.

The missing match is now more specific than an unnamed kernel bound.
The global consumers at Raw I, lines 112589--112604, require five
weighted output norms, including a graph norm containing
$\ell^3|\Delta_{\mathcal W}\delta A'|$ (lines 111408--111435).
The earlier response chart explicitly avoids applying covariant
Laplacians to general independent-current responses
(lines 15119--15141). A contraction in its weaker norm does not
establish that stronger output estimate. On the input side, the
field datum $\mathfrak j=\widetilde G_0(\sigma)\mathsf J$ is itself
assumed to satisfy a graph-norm bound at lines 72532--72535; treating
it as arbitrary weak-norm forcing would also lose a hypothesis.
The next proof must connect these actual input and output norms,
or justify a revised interface that controls the combined current
partner directly and still preserves every subsequent carrier domain.
Identities valid only when the current is the configuration's actual
Euler--Lagrange current cannot be used for an independent complex current.

There is also an upstream coupling-trajectory obligation. Raw I,
lines 56419--56441, makes its clause-zero statement conditional on
one-loop and unmarked-sector inputs. The concrete incomplete step at
lines 4072--4088 transfers an unmarked remainder estimate to the chosen
cutoff and radius; it is marked $(*)$. This is more precise than
treating every cited one-loop coefficient as unproved. The
effective-density representation is part of the joint induction.
Thus the coupling-flow inequalities used in the weighted calculation
above remain hypotheses for the construction. The original 1989 paper
also states its stability theorem under an effective-coupling
assumption; that historical statement does not by itself determine
what later work has proved.

\subsection{Two concrete repairs to test at the global interface}

The explicit earlier chart can at least be extended to two inputs,
without dropping the forcing. Freeze its background and decoupling
coefficients, write $G$ for its current-dependent Green operator, and
put $L_0=(I+G\Theta_J)H_0$. The equation
\[
 P=-\mathfrak j+L_0b-GV'(P)
\]
extends the block-only chart of Raw I, lines 14387--14429.
Assume $V'$ is holomorphic on a neighborhood of the radius-$2r$ ball,
$\norm G\le M_G$, $\norm{V'(P)}\le c\norm P^2$ there, and
$8M_Gcr\le1$. If
$d=\norm{\mathfrak j}+\norm{L_0}\norm b\le r/2$, the right side
maps the radius-$r$ ball into the radius-$5r/8$ ball and has Lipschitz
constant at most $1/2$. Indeed Cauchy gives
$\norm{V''(P)}\le4cr$ there. Thus its unique small solution is
holomorphic and $\norm P\le2d$.

For any weighted norm in which the same whole-ball Hessian bound is
$\norm{GV''(P)}_p\le1/2$, set
$R_P=(I+GV''(P))^{-1}$ and
$T_P=I-H^JD^{J\prime}(P)$. The layer of this earlier chart is
$\mathcal H=P-H^JD^J(P)$, and its exact data derivatives are
\[
 D_{\mathfrak j}\mathcal H=-T_PR_P,\qquad
 D_b\mathcal H=T_PR_PL_0.
\]
For $\mathfrak j=Gf$, the forcing derivative is
$D_f\mathcal H=-T_PR_PG$.
These are genuine two-input formulas, conditional on the stated
primitive and Hessian bounds; they do not identify this chart with
every later global layer.

Still taking $\mathfrak j=Gf$, at the undecoupled value $\sigma=1$ use
$G=(I-\bar G\Theta_J)^{-1}\bar G$ and the operator identities
$\bar G=P_0G_0$, $\Delta_a\bar G=P_0^T$ and
$\Delta_aH_0b=Q^*q_0\ell^{-1}b$ give the forced representation
\[
 \Delta_aP=P_0^T[-f+\Theta_JP-V'(P)]+Q^*q_0\ell^{-1}b.
\]
This controls the augmented elliptic operator without differentiating
$f$. It does not control the covariant Laplacian of
$P-H^JD^J(P)$ automatically. Nor does walk decoupling preserve an
inverse identity without an additional proof.

There is also an exact simplification of an actual global consumer.
Let $L_\sigma=\Delta_W-P_1(\sigma)+\mathsf K$ and retain the
deposited partner-current formula
\[
 \mathfrak s^\sharp_\sigma
   =L_\sigma A'-\mathfrak m_H(\sigma)X+F(W,\mathbb H).
\]
For a scalar cutoff $\chi$ whose arguments stay in the domain of $F$,
the cut current defined at Raw I, lines 111542--111543, satisfies
\[
 \begin{split}
 \mathcal J^\sharp[\chi]={}&\chi\mathfrak s^\sharp_\sigma
    +[L_\sigma,\chi]A'
    -[D_W^*D_W,\chi]H(\sigma)X\\
   &+F(W,\chi\mathbb H)-\chi F(W,\mathbb H).
 \end{split}
\]
This follows simply by expanding the first two terms; the separate
$\mathfrak m_H(\sigma)X$ terms cancel. Local second-order commutators
require first derivatives of the field and up to two derivatives of
$\chi$. For a nonlocal kernel $P_1$, its commutator has kernel
$P_1(x,y)(\chi(y)-\chi(x))$, so a weighted first-moment kernel bound
is the relevant input.

Consequently this cut consumer can potentially be controlled by an
\emph{independent bound on the combined current}, plus commutator
bounds, instead of separate Laplacian and $\mathfrak m_HX$ estimates.
The existing estimate of $\mathfrak s^\sharp$ uses those separate
estimates, so reversing it would be circular. The missing step is to
identify and bound a substituted current expression directly, and then
check all remaining consumers and carrier domains. The algebra above
is a concrete repair route, not that completed identification.

\subsection{Separate coefficient extraction from the cutoff estimate}

There is a useful logical refinement. The fourth hypothesis of Raw I's
unmarked-increment lemma, lines 4053--4057, already assumes the small
remainder estimate with its factor $\varepsilon_1^{\mathrm I}$.
The modified-cutoff theorem at lines 4170--4209 is the upstream claim
that must supply it for the actual fluctuation integral. The following
elementary extraction lemma shows what can be proved once that analytic
estimate and the correctly normalized response bounds are available.
It does not prove the change-of-cutoff estimate itself.

\begin{lemma}[Uniform quadratic extraction from pointed remainders]
Let $R_{X,z}(A)$ be holomorphic for $\norm A_X<r$, with $z\in X$
for finite nonempty $X\subset\mathbb Z^4$, $d(X)\ge0$, and
\[
 \sup_{\norm A_X<r}|R_{X,z}(A)|
       \le A_0\epsilon e^{-\kappa d(X)}.
\]
Assume, in the same normalized field norm,
\[
 \norm{h_x}_X\le H e^{-m\,\mathrm{dist}(x,X)},\quad
 \operatorname{diam}X\le M(d(X)+2),\quad
 \sup_z\sum_{X\ni z}e^{-\kappa_0d(X)}\le K_0,
 \qquad \kappa>\kappa_0.
\]
Here $r,m,M>0$, $\epsilon\ge0$, and all other bound constants are finite.
Put
\[
 a=\min\{m/2,(\kappa-\kappa_0)/(4M)\},\qquad
 C_* =4A_0H^2r^{-2}K_0e^{4aM},
 \quad S_j(a)=\sum_{u\in\mathbb Z^4}|u|^j e^{-a|u|}.
\]
Then the absolutely convergent kernel
\[
 \Pi(x,y,z)=\sum_{X\ni z}D^2R_{X,z}(0)[h_x,h_y]
\]
satisfies
\[
 |\Pi(x,y,z)|\le C_*\epsilon
       e^{-a(|x-z|+|y-z|)}.
\]
In particular, any quadratic coefficient with the displayed pointed
form obeys
\[
 \left|\sum_{u,z}u_\mu u_\nu\Pi(0,u,z)\right|
       \le4C_*S_0(a)S_2(a)\epsilon.
 \tag{2.3}
\]
All constants are independent of a scale index if the hypotheses are.
\end{lemma}
\begin{proof}
Apply two-variable Cauchy to $R_{X,z}(s h_x+t h_y)$ at radii
arbitrarily below $r/(2\norm{h_x}_X)$ and
$r/(2\norm{h_y}_X)$, then take their limits. The mixed derivative
is bounded by
\[
 4A_0\epsilon H^2r^{-2}e^{-\kappa d(X)}
 e^{-m(\mathrm{dist}(x,X)+\mathrm{dist}(y,X))}.
\]
Since $z\in X$,
$|x-z|+|y-z|\le\mathrm{dist}(x,X)+\mathrm{dist}(y,X)
+2\operatorname{diam}X$.
The chosen $a$ leaves at least $\kappa_0d(X)$ in the decay exponent;
the diameter allowance contributes $e^{4aM}$. Sum over $X$.
Finally write $u=z+(u-z)$ and use
$|u|^2\le2|z|^2+2|u-z|^2$. Summing the product of the two
exponentials gives $4S_0(a)S_2(a)$, proving (2.3).
\end{proof}

Thus coefficient summation is not the same missing assertion as
stability of the analytic remainder under the chosen cutoff. To apply
this lemma, the deposited remainder must have the pointed localization
and response/domain normalization just assumed, uniformly in the scale.
Those identifications must be checked explicitly. An ordinary global
supremum estimate or a real-field tail probability does not supply them.

\begin{lemma}[Correlated Gaussian cutoff cost]
Let $X_b\in\mathbb R^D$ be centered jointly Gaussian blocks. Suppose
the covariance of the joint vector $(X_b)_{b\in S}$ has operator norm
at most $\sigma^2$, $\sigma>0$, uniformly over the finite sets $S$ under consideration.
Then
\[
 \mathbb P\{\norm{X_b}_2>R\text{ for every }b\in S\}
 \le \left[2^{D/2}\exp\{-R^2/(4\sigma^2)\}\right]^{|S|}.
 \tag{2.4}
\]
No independence of the blocks is assumed. For
$R(g)=c(\log g^{-2})^p$, $c>0$, $p>1/2$, and fixed $C_1>0$, $q\in\mathbb R$,
the bracket in (2.4) is at most
$C_1g(\log g^{-2})^q$ for all sufficiently small $g$.
That threshold can be chosen independently of a later small parameter
$\epsilon$.
\end{lemma}
\begin{proof}
On the event in (2.4), $\sum_b\norm{X_b}_2^2>|S|R^2$.
The Gaussian exponential moment at parameter $1/(4\sigma^2)$ is
\[
 \det(I-\Sigma_S/(2\sigma^2))^{-1/2}\le2^{D|S|/2}.
\]
Markov's inequality proves (2.4), including singular covariances by
continuity. With $s=\log g^{-2}$, the logarithm of the ratio of its
bracket to $C_1g s^q$ is
\[
 \tfrac D2\log2-\log C_1-\frac{c^2}{4\sigma^2}s^{2p}
       +\tfrac s2-q\log s\longrightarrow-\infty.
\]
This proves the uniform eventual comparison, without reference to
$\epsilon$.
\end{proof}

Consequently the numerical conditions
$C_1g(\log g^{-2})^{p_1}\le\varepsilon_1^{\mathrm I}$ and a growing
logarithmic cutoff radius can pay the Gaussian cutoff factor in the
same order of parameter choices. What is still needed is the
\emph{localized analytic} remainder estimate on the fixed complex
background domain, including its normalization. A tail probability
for a real Gaussian cannot replace that estimate.

The original CMP119 discussion also narrows the required scope:
on printed p.\ 276, terms localized wholly inside the small-field
region retain the earlier conditions; the altered space (3.43) is
assigned to terms meeting its boundary \cite{balaban119}. The
whole-lattice unmarked calculation has no such boundary-crossing terms.
Its immediate obligation is therefore the changed real cutoff with
the fixed complex-background radii, rather than arbitrary layered
boundary domains.
Also, the deposited exponents obey $p_1<p_0$ (Raw I, lines
1355--1363). The auxiliary box $C_1(\log g^{-2})^{p_1}$ therefore
cannot automatically be identified with a domain containing the full
real cutoff of order $(\log g^{-2})^{p_0}$. If the actual comparison
needs $gR(g)\le\varepsilon_1^{\mathrm I}$, this can be imposed by a
further coupling ceiling chosen after $\varepsilon_1^{\mathrm I}$;
its sufficiency for the analytic argument must still be verified.

The preparatory-window relative boundary estimate in Raw II,
lines 2376--2385 and 2809--2812, remains a separate antecedent.
The contraction lemma does not prove it. Nor can real-source positivity
replace the complex estimate: for a Gaussian of variance
$\sigma_a^2\to\infty$, subtracting the exact holomorphic counterterm
$\sigma_a^2z^2/2$ makes the residual identically one, but the bound
$|Z(u+iv)|\leq Z(u)$ still loses $e^{\sigma_a^2v^2/2}$.

\section{A factorial bound without a shrinking global zero-free disc}

A many-cell bound $Q^n$ inserted into the preceding zero-free radius can
produce an artificial loss of order $Q^{n^2}$. The following exact
combinatorial argument avoids it.

\begin{lemma}[Coherent subset bound]
For each finite index set $I$, let $h_I$ be holomorphic on
$|z_i|\nu_i<r$, with $h_I(0)=1$ and
\[
h_I(z_B,0_{I\setminus B})=h_B(z_B),\qquad |h_B|\leq Q^{|B|}.
\]
The same $Q,r$ are used for every subset. Then
\[
|\partial_I\log h_I(0)|
\leq |I|!(2Q/r)^{|I|}\prod_{i\in I}\nu_i.
\tag{3.1}
\]
No zero-free polydisc of fixed radius is assumed.
\end{lemma}
\begin{proof}
Put $m_B=\partial_Bh_B(0)$. Cauchy gives
$|m_B|\leq(Q/r)^{|B|}\prod_{i\in B}\nu_i$.
The local logarithm at the origin satisfies
\[
\partial_I\log h_I(0)=
\sum_{\pi\in\mathcal P(I)}
(-1)^{|\pi|-1}(|\pi|-1)!\prod_{B\in\pi}m_B.
\]
If $S(n,k)$ is a Stirling number of the second kind, the sum of absolute
coefficients is $\sum_k(k-1)!S(n,k)$.
Counting ordered partitions by their positive block sizes gives
$k!S(n,k)\leq n!\binom{n-1}{k-1}$.
The coefficient sum is therefore at most $n!2^n$.
\end{proof}

The restriction identity requires a \emph{fixed} grid and a subtraction
functional that does not change when a source coefficient is set to zero.
For disjoint real supports,
\[
\norm{\Rea\sum_i z_i f_i}_\sharp+
\norm{\Ima\sum_i z_i f_i}_\sharp
\leq4\max_i|z_i|\nu_i.
\]
Thus an estimate for arbitrary finite real Lipschitz source families,
depending only on their combined source norm, can be applied pointwise
to the two profiles $\Rea w,\Ima w$. This gives a product source domain
with radius independent of $n$. The functional must be independent of
its presentation as a sum. The raw account advertises that interface;
its proof remains an obligation.

Centered moments also retain factorial growth: if
$|\cum_k|\leq k!A^k\prod\nu_i$ for every subset and $k\geq2$,
their exponential generating-function majorant is
$\exp((At)^2/(1-At))$. Cauchy at $t=(2A)^{-1}$ gives, for example,
$|S_n|\leq n!(4A)^n\prod\nu_i$.

\subsection{Keeping the cell side fixed}

The existing placement construction allows cell size to grow with $n$,
which loses the needed factorial control. The following repair keeps
one grid for every subset of a selected patch tuple; different tuples
may use different grids because the final cumulants are multilinear.
The underlying source estimates are still hypotheses.

\begin{theorem}[All-order flat-test bound]
Assume the fixed-scale source, denominator and locality estimates above,
with the following uniformities. An admissible fixed cell side $\ell$
can be used for the many-cell chessboard inequality; the RG grid and
reflection grid may be translated together on a mesh of size $O(a)$;
all constants are uniform in these translates and source positions.
The source counterterms are presentation-independent holomorphic
functionals on one fixed complex norm ball, vanish at zero, and obey
the anchored estimate (2.1) for unit balls, including reflected weights.
The bulk subtraction is kept fixed when source coefficients vanish.
Enlarged polymer supports that connect a buffered source to a face,
or connect separated sources, have the corresponding diameter lower
bound. Assume a common subsequential hierarchy of centered off-diagonal
moments has been extracted.

Then, on Schwartz tests flat to all orders on every pair diagonal,
this hierarchy extends uniquely and satisfies
\[
 |S_n(\Psi)|\leq C(n!)^{23}q_{15n}(\Psi),\qquad
 q_m(\Psi)=\max_{|\alpha|\leq m}\sup_X
       (1+|X|^2)^{m/2}|\partial^\alpha\Psi(X)|.
 \tag{3.2}
\]
Here $C$ is independent of $n$. If the anchored estimate is uniform
also for shrinking anchoring balls with its explicit $(R+\ell_j)^4$
factor, the factorial exponent improves from $23$ to $19$.
Neither exponent is asserted optimal.
\end{theorem}
\begin{proof}
\emph{Patch geometry.}
Choose a smooth spatial partition of unity on mesh
$\varepsilon_n=\min(1,\ell)/(64n)$ with patch supports of coordinate
width $2\varepsilon_n$. For one selected patch from each source,
at most $6n\varepsilon_n<\ell$ of the phases of an $\ell$-grid are
forbidden in each coordinate if a buffer $2\varepsilon_n$ is required.
Choose allowed phases and round to the admissible grid mesh. For
$a\leq c\varepsilon_n$, every patch retains buffer
$b_n\geq\varepsilon_n$. Keep this grid and its subtraction data for
all subsets of this tuple. There are at most $n$ occupied cells.
For a source supported in a ball of radius $R$,
\[
 \sum_s\norm{\chi_s^{\varepsilon_n}f}_{C^1}
 \leq C(1+R)^4 n^5\norm f_{C^1}.
 \tag{3.3}
\]
This follows from $O((1+R)^4\varepsilon_n^{-4})$ patches and the
first-derivative cost $O(\varepsilon_n^{-1})$.

\emph{Face cost and source radius.}
For a holomorphic functional $F$ on a norm ball of radius $r_*$,
$F(0)=0$ and $|F|\leq B$ imply
$|F(w)|\leq B\norm w/r_*$ by radial Schwarz.
Apply this to the absolutely summable face counterterms. It is
legitimate even for a conjugate-reflected weight: freeze that weight
and vary a complex scalar multiplying the entire weight, as allowed
by the norm-ball hypothesis. Unit-ball anchoring bounds the face cost
per patch by
\[
 J_n\leq C H_p(b_n)\leq Cn^4(1+\log n)^p.
\]
Let $C_0$ absorb the fixed comparisons between the complete complex
source norm, reflected weights, and $\max_i|z_i|\nu_i$.
Take $r_n=r_*/[C_0(1+2J_n)]$. On this smaller product domain the remaining
face costs are at most one per source. The fixed-cell cost is also
constant per occupied cell. Thus $|h_B|\leq Q^{|B|}$ with $Q$
independent of $n$, and (3.1) applies with $r=r_n$.

For the mixed derivative of the subtracted local functional, locality
kills every polymer not meeting all selected supports. Anchor at the
first patch and use its diameter at least their minimum separation
$\delta$ to obtain $C H_p(\delta)r_*^{-n}\prod_i\nu_i$.
A transported bulk term contributes only if all variables belong to
one cell; in that case use the same estimate in each reflected copy
and average. No volume factor remains.
Combining this with (3.1), (3.3), and multilinearity gives, for any
$\eta>0$,
\[
 |\kappa_n(f_1,\ldots,f_n)|
 \leq[C_\eta(1+R)^4]^n(n!)^{10+\eta}
       [1+H_p(\delta)]\prod_i\norm{f_i}_{C^1}.
 \tag{3.4}
\]
Here $n^{9n}n!$ supplies ten factorial powers and the logarithm fits
in $(n!)^\eta$ with an exponential constant.
The restriction $a\leq c/n$ is needed only eventually for each fixed
$n$, which suffices for a countable subsequence extraction.

\emph{Centered full moments.}
Only partitions with blocks of size at least two contribute. For $s\geq1$,
\[
 \sum_{\pi}\prod_{B\in\pi}(|B|!)^s
 \leq(n!)^{s-1}\sum_{\pi}\prod_{B\in\pi}|B|!
 \leq2^{n-1}(n!)^s.
\]
The last bound counts ordered partitions by positive compositions.
There are at most $n/2$ nonzero blocks and
$1+H_p(\delta)\leq C\delta^{-6}$ for $0<\delta\leq1$.
Hence, for supports in arbitrary translates of unit balls,
\[
 |S_n(f_1\otimes\cdots\otimes f_n)|
 \leq C_\eta^n(n!)^{10+\eta}\min(1,\delta)^{-3n}
                       \prod_i\norm{f_i}_{C^1}.
 \tag{3.5}
\]

\emph{Joint tests.}
This final step uses only a bound of the form (3.5) with factorial
exponent $s$. Put $\epsilon_j=2^{-j}$ and choose a fixed smooth
nondecreasing $\eta_0$ equal to zero on $[0,2]$ and one on $[4,\infty)$.
Set
\[
 F_j(X)=\prod_{a<b}\eta_0(|x_a-x_b|/\epsilon_j),\quad
 \chi_0=F_0,\quad\chi_j=F_j-F_{j-1}.
\]
These sum to one off the diagonals. On the support of $\chi_j$ or its
derivatives, all pairs are separated by $2\epsilon_j$ and, for $j\geq1$,
some pair is within $8\epsilon_j$.
Localize each variable on mesh $\rho_j=\epsilon_j/64$.
Enlarged one-variable chart supports are separated by at least
$\epsilon_j$. Fourier expansion on each rescaled product chart gives
tensor factors of $C^1$ norm at most
$C\rho_j^{-1}(1+|m_i|)$.
Applying $(1-\Delta_{y_i})^3$ for each variable bounds the sum of
the absolute Fourier coefficients times $\prod_i(1+|m_i|)$ by
$C^n$ times derivatives of order at most six in each variable.
The frequency sum converges because
$\sum_{m\in\mathbb Z^4}(1+|m|)(1+|m|^2)^{-3}<\infty$.

Each pair cutoff receives at most twelve derivatives; each spatial
bump receives at most six. The product rule therefore costs at most
$C^n n^{12n}$, using only fixed bounded-order cutoff derivatives.
For $j\geq1$, move one point of a close pair onto the other.
Taylor flatness of order $9n$ applied after at most $6n$ test
derivatives gives
\[
 |\partial^\alpha\Psi(X)|
 \leq C^n\epsilon_j^{9n}q_{15n}(\Psi)
                        \prod_i(1+|c_i|^2)^{-3},
 \qquad |\alpha|\leq6n,
\]
on the relevant chart. The segment moves only one coordinate a bounded
distance, so the product weight is bounded by
$C^n(1+|Y|^2)^{3n}$ along it and is absorbed by $q_{15n}$.
Chart positions sum with
$\sum_{c\in\rho_j\mathbb Z^4}(1+|c|^2)^{-3}\leq C\rho_j^{-4}$.
The net small-scale factor is therefore
\[
 \epsilon_j^{9n}\epsilon_j^{-3n}
 \epsilon_j^{-n}\epsilon_j^{-4n}=\epsilon_j^n.
\]
Its dyadic sum is bounded; $j=0$ needs no Taylor gain. Since
$n^{12n}\leq e^{12n}(n!)^{12}$, the result is
$C_\eta^n(n!)^{22+\eta}q_{15n}(\Psi)$.
Choose $\eta=1/2$ and absorb the exponential into another half factorial
power to obtain (3.2).

Spatial cutoffs followed by $F_J$ approximate a flat Schwartz test in
every Schwartz seminorm: arbitrary Taylor flatness beats all fixed
cutoff derivative losses. Thus compact tests away from diagonals are
dense, which proves existence and uniqueness of the stated extension.
No contact-term extension to all Schwartz tests has been chosen.

Finally, if shrinking-ball anchoring is uniform, each patch has radius
$R_n\asymp b_n\asymp n^{-1}$. Its face sum is bounded instead by
\[
 C\sum_{j\geq0}(1+j)^p(1+R_n/\ell_j)^4e^{-cb_n/\ell_j}
 \leq C(1+\log n)^{p+1}.
\]
Split at $\ell_j\asymp b_n$: the coarse sum has logarithmically many
terms and the fine sum decays exponentially in $L^j$.
This removes four factorial powers throughout, giving exponent $19$.
\end{proof}

The form of (3.2) is the linear-growth regularity condition used in OS
reconstruction; an explicit modern formulation appears in
\cite[Section 2.5, equation (22)]{amt}. That paper's two-dimensional
construction supplies no Yang--Mills input here.

Reflection positivity also extends, provided the hierarchy came from
one common centered sequence of reflection-positive lattice theories.
First approximate compact joint tests in the strict ordered positive-time
chamber by tensor sums supported in separated product boxes. Each finite
OS matrix is an entrywise limit of positive lattice matrices. Then use
the joint bound and density of interior compact tests in the Schwartz
space flat at the chamber walls. Reflected products are flat on all
collision diagonals, including the time-zero wall, so positivity
persists by continuity. This does not supply a second reflection,
full rotations, clustering, or nontriviality.

For the deposited one-frame torus geometry, a period $N=2L^k$ and
cell side $\ell=L^r$, $r\leq k$, give an even cell count
$N/\ell=2L^{k-r}$. Bare-step translations preserve the link-plane
faces and their antipodal partners. Thus odd $L$ is compatible with
this one-frame placement; no power-of-two cell count is required.
The remaining translation uniformity concerns the RG source estimates,
not this elementary tiling. Notebook 02 handles the additional
requirements of two-frame placement on the same odd-$L$ regulator.

There is also a concrete obstruction to removing the face buffer casually.
For a centered reflection-positive Gaussian pair $(X_+,X_-)$ with variances $A_a$
and covariance $B_a\geq0$, Cauchy--Schwarz yields
\[
|\E e^{zX_+}|
\leq \exp\!\left(\tfrac12 A_a\Rea z^2+\tfrac12B_a|z|^2\right).
\]
Subtracting the exact holomorphic counterterm $A_az^2/2$ leaves the
bound $e^{B_a|z|^2/2}$, even though the actual normalized function is
identically one. If $B_a$ diverges, a new estimate is needed.
The real part of a holomorphic function vanishing at zero cannot
uniformly cancel a positive divergent $|z|^2$ term: its circle average is
zero. Interface locality by itself therefore does not fix complex
chessboard discontinuities.

\section{Rotation restoration: a sharp conditional threshold}

The supplied rotation theorem applies to closed two-form laws with
Gaussian domination in an unweighted $L^2$ norm. It does not directly
apply to the composite scalar $\tr F^2$. In particular the non-Abelian
Bianchi identity is $d_AF=0$, rather than $dF=0$.

\subsection{A finite-lattice ingredient can be proved directly}

For the pure Wilson measure, diagonal reflection positivity itself has a
short proof on the correct finite-volume geometry.
Use the period lattice generated by
$n(e_0+e_1)$, $m(e_0-e_1)$ and transverse coordinate periods.
Reflection across $s=t+x=0$ fixes the sheets $s=0,n$ and separates
the two open halves. Only $(t,x)$ plaquettes cross a sheet.
Their two length-two half-path holonomies are $A$ and $B$, and their
Boltzmann factor is
\[
 e^{c\Rea\tr(AB^\dagger)},\qquad c\geq0.
\]
In real and imaginary matrix coordinates,
$\Rea\tr(AB^\dagger)$ is an ordinary real inner product.
Expanding its exponential in coordinate monomials gives a sum
$\sum_\alpha c_\alpha a^\alpha b^\alpha$ with $c_\alpha\geq0$.
Condition on links lying in the fixed sheets and integrate the remaining
links with product Haar measure. Each term is a squared half-space
integral. This proves reflection positivity for bounded half-observables,
including for $SU(2)$ and $SU(N)$.
Diagonal Wilson positivity is also used explicitly in \cite{kz}.

This does not settle \emph{same-limit} positivity. Ordinary coordinate
cubic tori have a different geometry. One must compare their selected
limit with the adapted construction, or choose a construction retaining
all required symmetries from the outset.
Centered free-boundary cubes retain both central reflection directions,
but would require a separate proof of translation invariance in their
limit. Changing boundary conditions does not remove that obligation.

\subsection{A compact-group mechanism for transferring insertions}

The same-limit argument in companion Notebook 02 exploits the compact
Wilson link variables. The following elementary lemma is its local
analytic input; the full proof additionally needs the geometry, slab
normalization, and limit argument given there.

\begin{lemma}[Wilson convolution conjugation]
On $G=SU(2)$ with normalized Haar measure and fixed $\beta>0$, let
\[
 k_\beta(g)=e^{(\beta/2)\tr g},\qquad
 (C_\beta\psi)(v)=\int_Gk_\beta(vu^{-1})\psi(u)\,du.
\]
If $f$ is a finite Peter--Weyl polynomial of largest spin $r$, then
$T_f=C_\beta^{-1}M_fC_\beta$, initially defined on finite representation
sums, extends continuously from $C^\infty(G)$ to itself. For each
Sobolev index $s$,
\[
 \norm{T_f\psi}_{H^s}\leq C_{\beta,f,s}
                         \norm\psi_{H^{s+2r}}.
\]
The reverse conjugation obeys the same type of estimate. Tensor-product
versions hold on any fixed finite set of links, with constants
independent of the ambient volume.
\end{lemma}
\begin{proof}
The spin-$j$ eigenvalue of $C_\beta$ is
$c_j=2I_{2j+1}(\beta)/\beta>0$. Indeed, integrate the character
$\sin((2j+1)\theta)/\sin\theta$ against
$(2/\pi)e^{\beta\cos\theta}\sin^2\theta\,d\theta$ and divide by
$2j+1$. Multiplication by $f$ connects only spins $k$ with
$|k-j|\leq r$. For integer $n\geq0$,
\[
 I_n(\beta)=\frac{(\beta/2)^n}{n!}S_n,\qquad
 1\leq S_n=\sum_{l\geq0}\frac{(\beta^2/4)^l n!}{l!(n+l)!}
               \leq e^{\beta^2/4}.
\]
Thus ratios of the eigenvalues on this finite band grow at most as
$C(1+j)^{2r}$. Each multiplication block has norm at most
$\norm f_\infty$, and only a bounded number of spin sectors are
connected. Weighted square summation proves the Sobolev bound;
Sobolev embedding gives continuity on $C^\infty$.
\end{proof}

In particular $C_\beta T_fB=fC_\beta B$. Since the kernel is symmetric,
this is the exact insertion identity
\[
 \int k_\beta(vu^{-1})f(u)B(v)\,dv
 =\int k_\beta(vu^{-1})(T_fB)(v)\,dv.
\]
On a fixed set of crossing links it can transfer a lower-layer
polynomial into a smooth function of the next layer. A local bounded
smooth perturbation $e^{tR}$ changes the resulting expression by
$O(t)$ in supremum norm, by the finite derivative loss and compactness.
These constants may depend on $\beta$; this is a fixed-lattice step.

Companion Notebook 02 completes the three additional interfaces at fixed
$\beta>0$: gauge-invariant layer removal using a conditional convolution
identity on the original link law; a clean-plane operator estimate from
exact physical transfer normalization; and two-frame placement using
alternating coordinate site/link walls. The latter retains the actual
side $n=2L^M$, $L$ odd, which cannot satisfy the convenient QED3
divisibility assumption. Diagonal feature kernels give the exact
cellwise folds. Full reflected source arrays are necessary in the first
Jensen comparison to retain a pressure error divided by the volume.
The assembled result is that every fixed-coupling coordinate-torus
subsequential infinite-volume limit has diagonal RP. Its proof takes
volume to infinity before removing the auxiliary tilt, uses convex
separation rather than uniqueness, and has been internally checked but
not externally certified. Passing this property to a renormalized
continuum hierarchy still requires that hierarchy's common convergence.

\subsection{A counterexample that survives the physical quotient}

Let $\phi$ be a massive free Euclidean scalar in $\R^4$ and put
$Y=P(\partial)\phi$, where $P(\partial)=\sum_{\mu=0}^3\partial_\mu^4$.
Then
\[
\widehat C_Y(p)=\frac{(\sum_\mu p_\mu^4)^2}{p^2+m^2}.
\]
This law has translations, signed coordinate permutations, factorial
Schwartz bounds, and reflection positivity at coordinate and $45$-degree
mirrors. Indeed $P$ is local and commutes with those reflections, so all
their positive-half-space observables pull back to those of $\phi$.

Nevertheless $Y$, declared to be a scalar, is anisotropic even off the
diagonal. For a $\pi/4$ rotation $R$, a contact-only covariance difference
would force $p^2+m^2$ to divide $P(Rp)^2-P(p)^2$.
At $p=(z,0,0,0)$ the latter equals $-3z^8/4$, nonzero at $z=im$.
This disproves divisibility. The OS quotient is nontrivial: its
one-particle multiplier is
$\omega(k)^4+\sum_jk_j^4\geq m^4$.
Thus broad regularity plus the two reflection directions is insufficient.

\subsection{The angular-degree lemma}

Let $M_N(T;\zeta)$ be an angular moment for an admissible $N$-tuple $T$.
Let $J$ denote antilinear time reflection. Assume:
\begin{enumerate}
\item $M_N$ is entire, $\pi/2$-periodic, and
\[
|M_N(T;\theta+i\beta)|\leq C_{N,T}e^{\kappa N|\beta|},
\quad 0\leq\kappa<4.
\]
Reality supplies both signs of every occurring frequency; $M_0=1$.
\item The tilted-frame moments obey
$\widetilde M_N(T;\zeta)=M_N(T;\zeta+\pi/4)$.
\item For every real $\beta$, both kernels
$M_{a+b}(J\chi,\eta;i\beta)$ and
$\widetilde M_{a+b}(J\chi,\eta;i\beta)$ are positive semidefinite on
positive-side tuples of \emph{all} lengths, including zero.
\item If $r_N>0$ is the largest absolute occurring frequency, the
coefficient at frequency $-r_N$ is nonzero on some tuple $(J\chi,\eta)$
for every split $N=a+b$, including an empty side.
\end{enumerate}
The last two requirements are analytic inputs, not consequences of
reflection positivity at real angle alone.

\begin{theorem}[Strict angular type below four]
Under these assumptions every $M_N$ is independent of angle.
\end{theorem}
\begin{proof}
Contour shifting kills Fourier frequencies $|\ell|>\kappa N$.
Quarter-turn periodicity allows only $\ell\in4\mathbb Z$.
Let $r_N$ be the largest absolute occurring frequency, zero if none.
Detection and Cauchy--Schwarz at imaginary angle give
\[
2r_{a+b}\leq r_{2a}+r_{2b}.
\]
At order $2a$, multiply the positive kernels by $e^{-r_{2a}\beta}$
and let $\beta\to+\infty$. Their leading coefficient kernels differ by
$e^{-ir_{2a}\pi/4}=(-1)^{r_{2a}/4}$. A nonzero kernel and its negative
cannot both be positive. Detection therefore gives $r_{2a}\in8\mathbb Z$.
Set $b_j=r_{2j}$. It is convex, $b_0=0$, and its increments are
nonnegative multiples of eight. Any positive increment forces eventual
slope at least eight, contrary to $b_j\leq2\kappa j<8j$.
Thus $r_{2j}=0$; the split $(0,N)$ gives $r_N=0$.
\end{proof}

The threshold is sharp for this abstract argument.
For the field $Y$ above, each insertion has angular frequencies
$0,\pm4$, and the two-point function has a nonzero eighth harmonic:
$(\cos^4\theta+\sin^4\theta)^2$ contains $\cos8\theta$.
Gaussian moments give maximal frequency $r_{2j}=8j$ and type $4N$.
The example also obeys the imaginary-angle kernel requirement:
the differential operator $P_{i\beta}$ commutes with antilinear
reflection, so free-field positivity applies to its transformed tests
in either frame. The leading coefficients are analytic off coincidence
and nonzero there; widely separated pairs retain their highest
harmonics by massive decay. Permutation and translation then place a
nonzero coefficient across each prescribed time split.

\subsection{A free-field benchmark for the desired exponent}

For the free four-dimensional Maxwell field, consider the separated
Schwinger kernels of its Wick-ordered scalar action density
$O=\sum_{\mu<\nu}:F_{\mu\nu}^2:$.
Split coordinates into a plane $(t,x)$ and two transverse coordinates $y$.
Smear each insertion in $y$ with $f_i\in C_c^\infty(\R^2)$.
For an ordered complex plane difference $z$ put
$q(z)=\Rea z_t-|\Ima z_x|>0$.
The partial transverse Fourier transform of each field-strength covariance
has the bound
\[
|\widehat K_{\alpha\beta}(z,k_\perp)|
\leq C\int_\R \omega\,e^{-q(z)\omega}\,dk_x,
\qquad \omega=(k_x^2+|k_\perp|^2)^{1/2}.
\]
The polynomial numerator from two derivatives is at most a constant
times $\omega^2$; the free spectral denominator is $2\omega$.
Scaling and Plancherel imply
\[
\norm{K_{\alpha\beta}(z,\cdot)}_{L^2(\R^2)}
\leq Cq(z)^{-3}.
\tag{4.1}
\]
For example, the momentum integral is bounded by
$Cq^{-2}(1+q|k_\perp|)e^{-q|k_\perp|/\sqrt2}$, whose
$L^2(d^2k_\perp)$ norm is $O(q^{-3})$.

Every Wick graph for these composite insertions is a union of cycles of
length at least two. The edge operator with kernel
$f_i(y)K_{ij}(y-y')$ has Hilbert--Schmidt norm
$\norm{f_i}_2\norm{K_{ij}}_2$.
The trace around each cycle is bounded by the product of these norms.
The cone distance is superadditive:
$q(z+w)\geq q(z)+q(w)$. Thus every nonadjacent ordered edge also has
distance at least $\delta$ when successive gaps do.
For $N\geq2$,
\[
|S_N(Z;f_1,\ldots,f_N)|
\leq C_N\delta^{-3N}\prod_i\norm{f_i}_2.
\tag{4.2}
\]
This proves the exponent-three benchmark in a free model.
It concerns separated kernels, not a claim that the coincident
Euclidean Wick square is an ordinary random tempered distribution.

For imaginary planar rotations,
$q(R_{-i\beta}z)\geq e^{-|\beta|}q(z)$.
If the \emph{interacting} ordered kernels satisfy the analogue of (4.2),
and the necessary analytic gluing, detection and two-frame Gram
identities are established, one obtains angular type $3N$.
Controlled logarithmic losses can be absorbed into $(3+\epsilon)N$,
with $\epsilon<1$.
That would meet the theorem's threshold.
Neither engineering dimension nor the free calculation proves this
interacting bound. Notebook 02 supplies same-limit diagonal RP at fixed
lattice coupling; its passage to the renormalized hierarchy and the
interacting analytic/Gram inputs remain necessary.

\section{A positive-energy witness at one coupling scale}

\begin{lemma}[Semigroup ratio]
Let $H\geq0$ be self-adjoint and $v\ne0$. If
$\langle v,e^{-tH}v\rangle\leq q\norm v^2$ with $0\leq q<1$, then for
$0<\alpha<\log(1/q)$,
\[
\norm{\mathbf1_{(\alpha/t,\infty)}(H)v}^2
\geq(1-qe^\alpha)\norm v^2.
\]
Moreover $(1-e^{-tH})v\ne0$, is orthogonal to every zero-energy vector,
and has norm at least $(1-q)\norm v$.
\end{lemma}
\begin{proof}
On $[0,\alpha/t]$, $e^{-tE}\geq e^{-\alpha}$. Integrate against the
spectral measure of $v$. The norm bound for $(1-e^{-tH})v$ follows from
Cauchy--Schwarz and its inner product with $v$.
\end{proof}

Here is a concrete test using the draft's own witness geometry. Set
\[
f_d(x)=\left(1-\frac{|x-(3d/4)e_0|^2}{(d/4)^2}\right)_+^2.
\]
It has supremum one, mass $\pi^2d^4/3072$, and Lipschitz constant
$32/(3\sqrt3\,d)<7/d$. Thus its test-class parameter can be
$\vartheta=10^{-3}$. Both pairs
$(\theta f_d,f_d)$ and $(\theta f_d,\tau_{4de_0}f_d)$ use the
\emph{same} scale $d$: their support separations are $d$ and $5d$.
Near-pair distances are at most $2d$, far-pair distances at least $5d$.
For the exact continuum kernel
$K(f,h)=\iint f(x)h(y)|x-y|^{-8}\,dx\,dy$,
\[
\frac{K(\theta f_d,\tau_{4de_0}f_d)}
 {K(\theta f_d,f_d)}
\leq(2/5)^8.
\]
The draft's two witness prefactors have ratio twelve. Consequently, if
its lower/upper witness and its sharp block-kernel asymptotic are proved,
\[
\limsup_{d\downarrow0}
\frac{\langle v_d,e^{-4dH}v_d\rangle}{\norm{v_d}^2}
\leq12(2/5)^8=0.00786432.
\tag{5.1}
\]
The running coupling cancels because the same $d$ is used.
One can then take $q=0.02$ for sufficiently small $d$.
With $\alpha=1$, more than $94.5\%$ of this vector's squared norm lies
above energy $1/(4d)$.

\subsection{Finite relative accuracy is enough}

In fact the sharp asymptotic in (5.1) is stronger than needed.
Let $K_-,K_+$ be the near and far continuum kernels, now with the common
normalization $6/\pi^4$. If their discrete background counterparts obey
\[
 (1-\epsilon)K_\pm\leq K_{{\rm bg},\pm}
             \leq(1+\epsilon)K_\pm,
\]
the same witness gives
\[
 \frac{\langle v_d,e^{-4dH}v_d\rangle}{\norm{v_d}^2}
 \leq12\frac{1+\epsilon}{1-\epsilon}(2/5)^8.
\]
Already $\epsilon=1/3$ gives $0.01572864<0.02$.
Thus fixed relative accuracy on two shapes suffices for the stated
positive-energy conclusion. The draft's coarse kernel constants alone
do not supply this accuracy.

Here is an independent free-lattice calculation and the exact
additional estimate that would supply it. In midpoint conventions put
$s_\mu(q)=2\sin(q_\mu/2)$ and $\lambda(q)=\sum_\mu s_\mu(q)^2$.
The covariance multiplier of one colour of the exact two-form Gaussian is
\[
 m_{\mu\nu,\rho\sigma}(q)=
 \frac{\delta_{\nu\sigma}s_\mu s_\rho-
 \delta_{\nu\rho}s_\mu s_\sigma-
 \delta_{\mu\sigma}s_\nu s_\rho+
 \delta_{\mu\rho}s_\nu s_\sigma}{\lambda(q)}.
\]
Write $m_0$ for the homogeneous multiplier obtained by replacing $s$ by
$q$, and $D_{PQ}(r)$ for its inverse Fourier transform away from zero.
The potential covariance is normalized to
$\delta_{\mu\nu}/(4\pi^2|r|^2)$.

\begin{proposition}[Free off-diagonal approximation]
For bare spacing $a$ and plaquette-midpoint separation $r$ with
$|r|\geq C_0a$, the infinite-lattice covariance satisfies
\[
 |P_{{\rm ex},a}(p,q)-a^4D_{\pi(p)\pi(q)}(r)|
 \leq Ca^6|r|^{-6}.
\]
\end{proposition}
\begin{proof}
Near zero,
$|\partial^\alpha(m-m_0)(q)|\leq C_\alpha|q|^{2-|\alpha|}$.
A smooth dyadic decomposition into $|q|\asymp2^{-j}$ therefore gives
Fourier pieces bounded by
$C_M2^{-6j}(1+2^{-j}|z|)^{-M}$ for any fixed $M>6$.
Summing on both sides of $2^{-j}|z|=1$ gives $C|z|^{-6}$.
The lattice multiplier away from zero is smooth and periodic in the
basepoint convention and has rapidly decreasing Fourier coefficients.
For the continuum multiplier, sufficiently high derivatives of its
high-frequency part are integrable; integration by parts gives the
same rapid decrease away from contact terms.

For different orientations, apply the same midpoint phase
$e^{iq\cdot(\delta_P-\delta_Q)}$ to both multipliers. Their difference
still has order two at zero, and its integer Fourier coefficient at
$n$ compares kernels at $z=n+\delta_P-\delta_Q=r/a$.
Finally use $D(r/a)=a^4D(r)$.
\end{proof}

Suppose the actual background covariance has been identified as
$C_{{\rm bg},a,h}=P_{{\rm ex},a}-C_{{\rm fl},a,h}$.
The additional input needed is the uniformly normalized
\emph{curvature} bound
\[
 |C_{{\rm fl},a,h}(p,q)|
 \leq Aa^4h^{-4}e^{-c|r|/h},\qquad |r|\geq h.
 \tag{5.2}
\]
Here $h$ is the physical witness-cell scale. A bound on the potential
propagator alone does not give (5.2) without its discrete curl estimates.
A scale decomposition with terms bounded by
$A_0a^4h_j^{-4}e^{-c_0|r|/h_j}$, $h_j=hL^{-j}$, would suffice because
\[
 \sum_{j\geq0}L^{4j}e^{-c_0\rho L^j}
 \leq C_{L,c_0}e^{-c_0\rho/2},\qquad \rho\geq1.
\]
The elementary summation is proved by factoring out the right-hand
exponential and using the summable majorant
$\sum_jL^{4j}e^{-c_0L^j/2}$.

For clarity, the discrete quadratic kernel meant here is
\[
 K_{{\rm bg},\pm}=\sum_{p,q}f_\pm(x_p)g_\pm(x_q)
                       |C_{{\rm bg},a,h}(p,q)|^2,
\]
with each positive plaquette orientation occurring once. The two
covariance factors already supply the $a^8$ Riemann-sum weight.
Any common colour or energy normalization must multiply both the
continuum and discrete kernels.

\begin{theorem}[Two-scale criterion]
Under the covariance identification and (5.2), the two fixed shapes
satisfy the relative comparison with
\[
 \epsilon\leq C_*\left\{a/d+(d/h)^4e^{-c\,d/h}\right\},
 \tag{5.3}
\]
for sufficiently small $a/d$, with a fixed reduction in $c$ if needed.
In particular finite thresholds $a/d\leq\eta_*$ and $d/h\geq R_*$
give $\epsilon\leq1/3$.
\end{theorem}
\begin{proof}
All distances lie in $[d,6d]$. The free approximation contributes
$O((a/d)^2)$ against the curvature scale $a^4d^{-4}$, and (5.2)
contributes $O((d/h)^4e^{-cd/h})$. Moving orientation midpoints to a
common basepoint costs $O(a/d)$. For the six by six curvature matrix,
\[
 \norm{D(r)}_{\rm F}^2=\frac6{\pi^4|r|^8},\qquad
 |\norm{D+E}_{\rm F}^2-\norm D_{\rm F}^2|
 \leq2\norm D_{\rm F}\norm E_{\rm F}+\norm E_{\rm F}^2.
\]
Riemann sums for the fixed Lipschitz shapes cost $O(a/d)$ after dilation
by $d$. Divide the absolute errors by the two fixed positive continuum
kernels and absorb quadratic errors, reducing the exponential rate if
necessary. This proves (5.3).
\end{proof}

The free finite-torus covariance tends to its infinite-volume value at
fixed $a$ by Fourier Riemann sums. This supplies no uniform torus error
as $a\to0$. Take volume first, or choose a sufficiently large volume
diagonally for these two probes, and separately verify the conditioned
covariance's volume limit. The deposited choice of $h$ has $d/h$
growing polylogarithmically as $d\downarrow0$; it would provide the
required separation once (5.2) is proved for that exact block map.

The interacting lower/upper witness and a valid reconstruction remain
required. The result proves no mass gap: arbitrarily small positive
energies in other states remain possible.

\section{Why this still would not identify Yang--Mills}

There is an explicit generalized-free countermodel to the logarithmic
two-point signature. In four dimensions use positive spectral density
\[
\rho(s)=\frac{a s^2}{\log^2(e+s/\Lambda^2)},\qquad
C(r)=\int_0^\infty \rho(s)
\frac{\sqrt s}{4\pi^2r}K_1(r\sqrt s)\,ds.
\]
It defines a local generalized free Wightman field with positive energy,
unique vacuum and clustering. Substitution $u=r\sqrt s$ and
$\int_0^\infty u^6K_1(u)\,du=384$ give
\[
C(r)\sim\frac{48a}{\pi^2}
\frac{r^{-8}}{\log^2(1/(\Lambda r))}.
\]
Thus this behavior does not by itself characterize non-Abelian dynamics.
The general spectral construction is standard \cite{dr}; this choice
and the displayed asymptotic provide the explicit counterexample.

Even a nonzero fourth connected function does not repair that criterion.
Take a generalized free $\chi$ with
$\rho_0(s)=a/\log(e+s/\Lambda^2)$ and its local Wightman Wick square
$O=:\chi^2:$. At separated Euclidean points,
\[
C_\chi(r)\sim\frac{a}{2\pi^2}
\frac{r^{-4}}{\log(1/(\Lambda r))},\qquad
C_O(r)=2C_\chi(r)^2
\sim\frac{a^2}{2\pi^4}
\frac{r^{-8}}{\log^2(1/(\Lambda r))}.
\]
The connected four-point function of $O$ is sixteen times the sum of
the three distinct four-cycles of $C_\chi$, and is nonzero.
The forward-cone spectral measure has polynomial growth; Wick-square
smearing is controlled by its positive-cone convolutions. A full
operator-domain presentation would accompany any publication use.
These are Wightman/separated-Euclidean examples, not unsubtracted
Euclidean probability covariances at contact.

There is also a separate vacuum issue. Adding an independent constant
random mode preserves the ultraviolet signature while creating
degenerate zero-energy states. The temporal difference in Section 5
removes that contamination.

If an actual Euclidean probability measure with a local bounded
observable algebra is available, phase selection has a useful route:
translation-invariant events are measurable from either time half-space
by translated Følner averages. Reflection positivity forces reflection
to act trivially on that invariant algebra. Applying positivity to
$\sqrt h A$, with $h$ invariant and nonnegative, then gives positivity
in almost every translation-ergodic component. After full Euclidean
covariance, these components cluster; this follows from the relative
Howe--Moore property of $(E(d),\R^d)$ for $d\geq2$ \cite{hm}.
The separated composite functions currently under discussion do not
automatically provide that measure/algebra. An algebraic phase-selection
argument or a direct clustering proof is still needed in their actual
framework.

\section{The next research sequence}

\begin{enumerate}
\item \textbf{Close one sourced induction at a fixed stopping scale.}
Use the joint fluctuation bounds of Notebook 17, the overlapping
integration theorem of Notebook 18, and the physical Gaussian and
measure factors of Notebook 19. Notebook 20 closes their coefficient
and precision match for the first bare Wilson small-field integral.
Notebook 21 extends the local decomposition and one classical-reference
integration uniformly through the minimized classical actions.
Notebook 22 proves stationary-background substitution with explicit
radius and support-weight losses. Notebook 23 proves local gauge
invariance of the complete correction and its quadratic Ward
identity. Notebook 24 gives the conditional inherited-interaction
bound and an explicit sufficient next-step norm test. Notebook 26
satisfies that test with a persistent auxiliary-cell convention,
returning one fixed positive support exponent for both inherited
and new terms. Notebook 27 supplies an actual composite macrostep,
including its exact local integral. Notebook 28 uses the smaller
representative of Notebook 25 to return the same positive analytic
radius at one step. It also contracts coefficientwise fifth-order
background remainders and the quantum response to differences
with matching low jets, at a sufficiently large fixed macro ratio
and small coupling. Notebook 29 proves the actual local spatial
symmetries, the total nonlinear Ward completion of the marginal
term, and a connected localization of canceled low kernels with
gain \(C(J)L^{-\alpha}\) under its coherent-family and exact
infinite-kernel hypotheses. Its classical semigroup permits
the finite endpoint ratio \(L=R^q\) to be chosen after fixing
the elementary \(R,J\). Its joint compactness and complete
interaction quotient justify the classical change of representation
under the stated map hypotheses. Notebook 30 establishes the
universal classical gauge identity, canonical periodic descent
and local continuity of the actual homogeneous quantum update.
Its varying-input theorem identifies finite histories under
their explicit intermediate-family budgets. Exact classical
reference transport permits a projected finite-block contraction
and an amplitude envelope on the selected infinite quotient
while the couplings remain weak. Notebook 31 proves the
corresponding finite local induction: finite Ward repair controls
wrapping terms, and exact equivalent representatives meet every
intermediate family budget. It also proves fixed-threshold
coupling parity and exact preservation of the infinite Maxwell
coefficient by classical substitution. On the infinite amplitude
envelope the drift is the pivot/determinant coefficient plus a
uniform error tending to zero with the coupling. The leading
coefficient must still be computed. Notebook 32 gives qualitative
uniform finite leakage, real-coupling continuity, an explicit
convergent determinant trace, and an exact leading determinant
cocycle. Its gauge/coarea identity supplies the actual coupled
vector, ghost and coarse comparison terms, and bounds the shifted
coarse effective-Hessian contribution uniformly in depth.
A positive block margin
in every varying-depth structural limit would transfer to finite
regulators without a quantitative rate.
Under a positive averaged drift hypothesis for the actual
finite leading coefficients, its shooting argument produces
arbitrarily long finite weak local trajectories with a fixed
endpoint. Establish that drift through an actual uniformly
controlled determinant comparison; it is not supplied by absolute
increment bounds or the cocycle.
Notebook 33 gives an exact constant-Cartan-curvature expression
for the actual Maxwell coefficient and removes its multiplier
curvature through second order using the classical envelope.
Its diagonal trace densities have individual gauge identities,
which are essential for matching the affine coefficient to the
total Hessian moment. The notebook computes a specified massive
continuum reference and rewrites the actual comparison as a
heat integral. Uniform bounds on the short-time, intermediate
and long-time parts remain to be proved, with the intermediate
coefficient matched to the continuum value and the soft
based-ghost modes controlled in the combined expression.
The based-ghost short-time part alone now has a uniform
bound from a closed-walk expansion on the actual pinned domain.
Notebook 34 extends this to the entire based-ghost coefficient
and proves its additional logarithmic interval from \(K^2\)
to \(K^4\). It also proves the full vector ultraviolet bound
using a uniform counting-\(\ell^2\) bound on the actual
composite constraint and a localized trace on fixed coarse
fibers. The short-time interval is therefore settled.
The remaining vector estimate must match its intermediate
coefficient and cancel the ghost's additional long-time
logarithm, or supply an equivalent bounded signed comparison.
The same notebook gives an exact finite physical quotient,
including both orbit-density determinants, and proves its
uniform flat \(cK^{-2}\) gap. This supplies an alternative
route to the combined comparison. Establish the joint
infrared limit and the required curvature-source bounds
before using the separate projected factors as infinite
affine trace densities.
Notebook 35 supplies that fixed-endpoint infrared bridge by
changing the exact tangent gauge to covariant block means.
It proves an actual local scalar/vector determinant identity,
joint analytic continuation through flat stabilizers, individual
gauge-invariant diagonal densities, and physical-scale flat
gaps with uniform heat and inverse-power estimates for both
new factors. Use this pair for the next large-endpoint curvature
comparison. Its scalar projector, corrected tangent constraint,
and full residual Hessian must be retained. The old point-pinned
ghost logarithm cannot be assigned to the new scalar factor.
The constant-flat anomaly found in the old physical quotient
has canceled algebraically before the new local limits.
Notebook 36 further removes the hard block projections by exact
scalar and vector Schur factorizations. It proves that the two
coarse determinant corrections have uniformly bounded actual
Maxwell coefficients. Its physical relative-vertex theorem
retains the true effective Hessian, using local path gauges
and the ordered factorization of the flat kernel; the
component-skew lattice terms are not discarded. Continue with
the resulting fine soft pair and compare its marginal
vertices and infrared penalties to a specified reference.
The existing Ward extraction matches momentum jets through
field degree four, but the required determinant remainder
must still be bounded uniformly in preceding depth.
The same notebook bounds the soft pair's ultraviolet and
infrared ends, so it suffices to prove the actual intermediate
heat coefficient \(11\alpha_g\) with an error integrable in
logarithmic time. Flat heat bounds alone do not identify this
coefficient.
Match the continued local representation
to the measure, and control the complement of the small-field chart.
Fix $h$, an endpoint inverse-coupling interval, the vacuum normalization,
and the exact class of reflected source profiles. Prove the denominator
window, the residual complex-source upper bound, and the anchored
counterterm sum (2.1). Track every constant in $a$ and $|\Lambda|$.
The four-stage sufficient criterion in Section 2 is only an auxiliary
tool. Audit the selected global two-input layer, its actual carrier
domains and inheritance/arrest closure, together with the coupling
trajectory and preparatory-window inputs. Reconcile the three-link
localized description with that global route before treating its
unresolved assumptions as prerequisites.
\item \textbf{Construct separated limits and preserve positivity.}
Use the fixed-cell theorem and diagonal extraction. Prove the
distance-uniform per-ball strengthening separately. Do not undertake
two-cutoff uniqueness as a prerequisite.
\item \textbf{Make the all-order bounds usable.}
Verify the translated-grid and presentation-independent holomorphy
hypotheses of Section 3. Its patch argument, factorial lemma and
joint-test extension then give the explicit linear-growth estimate
(3.2). The remaining reconstruction assumptions, including symmetries
and the vacuum condition, must be checked separately.
\item \textbf{Prove the interacting planar insertion estimate.}
Use the free exponent-three calculation as a benchmark, not as a
substitute. Use the same-state lattice positivity of Notebook 02, prove
its passage to the renormalized hierarchy, and establish the
analytic/Gram inputs of Section 4. The existing rotation-current route
is an alternative only if its far-boundary flux condition is proved;
renaming it does not establish it.
\item \textbf{Prove the two-shape kernel ratio.}
Verify the constrained fluctuation-curvature bound (5.2), its volume
limit, and the interacting witness. The free calculation and (5.3)
then give one-third relative accuracy, already enough for
positive-energy nontriviality through Section 5.
\item \textbf{Extend and identify the observable theory.}
Introduce scalar, tensor and other gauge-invariant sources, control
their operator mixing and joint limits, and match the required
Yang--Mills short-distance structure. Prove a suitable vacuum-sector
statement. Keep a mass-gap theorem as a distinct infrared task.
\end{enumerate}

\noindent\textbf{What this notebook has achieved.}
It reduces the first construction to explicit interfaces, proves several
conditional steps and tests them against counterexamples. The largest
unpaid debt remains the uniform \emph{sourced} non-Abelian RG analysis.
None of the deductions removes that debt.

\section*{Audit locations in the deposited account}
\small
These are navigation references to the pinned version, not assertions
that every dependency at those locations has been verified.
\begin{itemize}
\item \texttt{theorem\_0\_complete\_account\_raw\_III.tex},
lines 110862--110965: public proof-status ledger.
\item Raw II, lines 60700--60765: local source terms, analyticity,
stability and denominator floor; literal finite-family source domain.
\item Raw II, lines 2791--2979 and 59645--59855: source correlation
theorem and its final-scale consumer; lines 15206--15331 and
25412--25480 onward: sourced large-field step and joint induction.
\item Raw I, lines 73107--73123 and 73545--73634: the selected global
four-layer route and exclusion of the alternative localized chain;
lines 56419--56441: conditional running-coupling input.
The three-link discussion at Raw II, lines 48058--48320, is a distinct
description of site (L); its relation to the active global-carrier route
requires reconciliation before it identifies a necessary source gap.
\item Raw I, lines 155039--155083:
uniform characteristic-window proposal.
\item Raw II, lines 63365--63367:
outstanding final-scale statement; lines 63964--64274:
per-ball recentering/chessboard reduction.
\item Raw II, lines 145173--145222:
outstanding source lemma and continuum-kernel proof.
\item Raw III, lines 37365--37427:
weak fine/coarse propagator stationarity.
\item \texttt{ym4\_manuscript\_zeroth\_draft.pdf},
Section 4.3, pp.\ 68--75: flat-diagonal test spaces and all-order
growth issue; Definition 5.1 and Remark 5.7, pp.\ 77--78:
witness geometry and claimed block-kernel asymptotic.
\item \texttt{QED3\_slop/YM\_Rotation\_Invariance\_simplified.pdf}:
Theorem A, p.\ 2; angular and two-frame argument, pp.\ 7--12;
insertion bounds, pp.\ 13--16; contact lemma, pp.\ 17--18.
\end{itemize}
\normalsize

\begin{thebibliography}{9}
\bibitem{repo} Public experiment repository, pinned commit above.
\href{https://github.com/ethanwsussman/Yang--Mills-LLM-experiment}
{Repository and source documents}.
\bibitem{clay} A. Jaffe and E. Witten,
\emph{Quantum Yang--Mills Theory}, especially pp.\ 5--6.
\href{https://www.claymath.org/wp-content/uploads/2022/06/yangmills.pdf}
{Official problem statement}.
\bibitem{os2} K. Osterwalder and R. Schrader,
\emph{Axioms for Euclidean Green's Functions II},
Commun.\ Math.\ Phys.\ 42 (1975), 281--305.
\href{https://doi.org/10.1007/BF01608978}{Publisher record}.
The corrected reconstruction theorem's precise hypotheses must be
checked in any final construction.
\bibitem{amt} M. S. Adamo, Y. Moriwaki and Y. Tanimoto,
\emph{Osterwalder--Schrader axioms for unitary full vertex operator
algebras}, Section 2.5, equation (22), and Section 3.2.
\href{https://arxiv.org/abs/2407.18222}{arXiv:2407.18222}.
Used only for the explicit linear-growth formulation, not as a
four-dimensional existence theorem.
\bibitem{balaban122b} T. Ba\l aban,
\emph{Large field renormalization. II. Localization, exponentiation,
and bounds for the R operation}, Commun.\ Math.\ Phys.\ 122 (1989),
355--392. \href{https://doi.org/10.1007/BF01238433}{Publisher record};
\href{https://scispace.com/pdf/large-field-renormalization-ii-localization-exponentiation-h8b7ih93xo.pdf}
{accessible copy of the original paper}.
\bibitem{balaban119} T. Ba\l aban,
\emph{Convergent renormalization expansions for lattice gauge theories},
Commun.\ Math.\ Phys.\ 119 (1988), 243--285, especially p.\ 276,
the distinction preceding (3.43).
\href{https://scispace.com/pdf/convergent-renormalization-expansions-for-lattice-gauge-29jyexd38e.pdf}
{Accessible copy of the original paper}.
\bibitem{seiler} E. Seiler, \emph{Uses of Reflection Positivity},
especially pp.\ 4--5 and 15--17.
\href{https://ethz.ch/content/dam/ethz/special-interest/phys/theoretical-physics/itp-dam/documents/graf/schrader/Seiler.pdf}
{Author's presentation}.
\bibitem{dr} M. D\"utsch and K.-H. Rehren,
\emph{Generalized free fields and the AdS--CFT correspondence}.
\href{https://arxiv.org/abs/math-ph/0209035}{arXiv:math-ph/0209035}.
\bibitem{hm} R. Cluckers, Y. de Cornulier, N. Louvet, R. Tessera
and A. Valette, \emph{The Howe--Moore property for real and
$p$-adic groups}, Theorem 5.2 and Proposition 2.2.
\href{https://rcluckers.perso.math.cnrs.fr/prints/CCLTV.pdf}
{Author-hosted paper}.
\bibitem{kz} V. Kazakov and Z. Zheng,
\emph{Bootstrap for Lattice Yang--Mills theory}, p.\ 3 and Fig.\ 5.
\href{https://arxiv.org/abs/2203.11360}{arXiv:2203.11360}.
\end{thebibliography}
\end{document}
